\documentclass[a4paper,10pt]{report}\input{../0.Préambule/Préambule_cours_prof}\begin{document}

\newpage
\section*{Les parallélogrammes}
\addcontentsline{toc}{section}{Les parallélogrammes}

\subsection*{Propriétés des parallélogrammes}
\addcontentsline{toc}{subsection}{Propriétés des parallélogrammes}

\begin{exer1}

\begin{flushright}
\includegraphics[scale=1]{../40.Image/13. Exo-Les parallélogrammes 1.png}
\end{flushright}

\vspace{-2.5cm}
\begin{enumerate}[font=\color{exer1}]
\item Écris tous les noms possibles du parallélogramme ci-contre.
\item Sur la figure, repasse en vert le coté opposé à $[PA]$, en bleu un coté \\ consécutif à $[PA]$, en rouge l'angle opposé à $\widehat{PCR}$ et en violet un  \\ angle consécutif à $\widehat{PAR}$.
\item Complète les phrases suivantes avec l'un ou plusieurs des mots : \begin{center}
consécutifs \quad \quad diagonales \quad \quad opposés \quad \quad cotés  \quad \quad  angles.
\end{center}

Dans le parallélogramme $PARC$ :
\begin{enumerate}[font=\color{exer1}]
\item Les \dotfill $[PA]$ et $[CR]$ sont de même mesure.
\item La somme des deux \dotfill $\widehat{PAR}$ et $\widehat{ARC}$ fait $180$\degre .
\item Les \dotfill $[PR]$ et $[CA]$ se coupent en leur milieu.
\item Les \dotfill $\widehat{CPA}$ et $\widehat{CRA}$ sont de même mesure.
\item Les \dotfill $[PC]$ et $[RA]$ sont parallèles.
\end{enumerate}
\end{enumerate}
\end{exer1}

\begin{exer1}
Dans la figure ci dessous, les droites d'une \\ même couleur sont parallèles.

\vspace{-1.2cm}
\begin{flushright}
\includegraphics[scale=1]{../40.Image/13. Exo-Les parallélogrammes 2.png}
\end{flushright}

\vspace{-2.8cm}
\begin{enumerate}[font=\color{exer1}]
\item Nomme tous les parallélogrammes de cette figure.
\item Pourquoi peux-tu affirmer que ce sont des \\ parallélogrammes ?
\end{enumerate}
\end{exer1}

\vspace{0.5cm}
\begin{exer1}
Dessine et code le parallélogramme $ABCD$ selon les consignes et cite dans chaque cas la propriété du parallélogramme qui t'a permis de réaliser la consigne.

\begin{center}
\renewcommand{\arraystretch}{1.8}
\begin{tabularx}{\linewidth}{|>{\centering}m{4cm}|>{\centering}m{4.5cm}|>{\centering \arraybackslash}X|}
\hline
\rowcolor{exer1!20}Figure & Consigne & Propriété\\\hline
 & Code les cotés de même longueur & \\\cline{2-3}
 & Colorie d'une même couleur les angles de même mesure & \\\cline{2-3}
 & Code les longueurs égales sur les diagonales & \\\hline
\end{tabularx}
\end{center}
\end{exer1}

\noindent \begin{minipage}[t]{8cm}
\begin{exer1}
Complète les étiquettes sachant que $ROSE$ est un parallélogramme. Justifie tes réponses.

\begin{center}
\includegraphics[scale=1.2]{../40.Image/13. Exo-Les parallélogrammes 3.png}
\end{center}
\end{exer1}
\end{minipage}
\hspace{5mm}
\begin{minipage}[t]{8.5cm}
\begin{exer1}
La figure est dessiné a main levée

\begin{center}
\includegraphics[scale=0.8]{../40.Image/13. Exo-Les parallélogrammes 4.png}
\end{center}

\begin{enumerate}[font=\color{exer1}]
\item Complète les étiquettes sachant que $BLEU$ est un parallélogramme.
\item Justifie ta réponse pour l'angle $\widehat{BLE}$
\item Justifie ta réponse pour la longueur $BU$
\end{enumerate}
\end{exer1}
\end{minipage}

\vspace{-6mm}
\noindent \begin{minipage}[t]{8cm}
\begin{exer1}
On considère le parallélogramme $ABCD$

\vspace{-3mm}
\begin{center}
\includegraphics[scale=1]{../40.Image/13. Exo-Les parallélogrammes 5.png}
\end{center}

\begin{enumerate}[font=\color{exer1}]
\item Quel est la mesure de l'angle $\widehat{CBA}$ ?
\item Pourquoi ?
\end{enumerate}
\end{exer1}
\end{minipage}
\hspace{2mm}
\begin{minipage}[t]{8.5cm}
\begin{exer1}
On considère le parallélogramme $UVWT$

\vspace{-3mm}
\begin{center}
\includegraphics[scale=1]{../40.Image/13. Exo-Les parallélogrammes 6.png}
\end{center}

\begin{enumerate}[font=\color{exer1}]
\item Quel est la longueur $TW$ ?
\item Pourquoi ?
\end{enumerate}
\end{exer1}
\end{minipage}

\begin{exer1}
$EFGH$ est un parallélogramme.

\begin{enumerate}[font=\color{exer1}]
\item Justifie que $(EF)$ est parallèle à $(GH)$ et que $(EH)$ est parallèle à $(FG)$
\item Quelle est la somme des quatre angles de ce parallélogramme ? Justifie.
\end{enumerate}
\end{exer1}

\begin{exer1}
$ABCD$ est un parallélogramme de centre $O$. Justifie que $O$ est le milieu du segment $[AD]$
\end{exer1}

\subsection*{Propriétés des parallélogrammes particuliers}
\addcontentsline{toc}{subsection}{Propriétés des parallélogrammes particuliers}

\begin{exer1}
Code les longueurs égales et les angles droits, sachant que le quadrilatère est :

\begin{enumerate}[font=\color{exer1}]
\begin{tabularx}{\linewidth}{*{4}X}
\item un rectangle &
\item un losange &
\item un carré &
\item Un parallélogramme\\
\end{tabularx}
\end{enumerate}

\includegraphics[scale=0.95]{../40.Image/13. Exo-Les parallélogrammes 7.png}
\end{exer1}

\noindent \begin{minipage}[t]{10cm}
\begin{exer1}
Sans justifier, complète les étiquettes sachant que $EFGH$ est un losange et $KLMN$ est un carré tel que $KM = 7$cm.

\includegraphics[scale=0.95]{../40.Image/13. Exo-Les parallélogrammes 8.png}
\end{exer1}
\end{minipage}
\hspace{2mm}
\begin{minipage}[t]{6.5cm}
\begin{exer1}
On considère un carré $KLMN$ de centre $S$ et tel que $SM=2,7$cm

\begin{enumerate}[font=\color{exer1}]
\item Fais une figure à main levée ?
\item Quelle est la mesure de l'angle $\widehat{NSM}$ ? Pourquoi ?
\item Quelle est la longueur $NS$ ? Pourquoi ?
\end{enumerate}
\end{exer1}
\end{minipage}

\noindent \begin{minipage}[t]{8cm}
\begin{exer1}
On considère le rectangle $ABCD$.

\begin{center}
\includegraphics[scale=1]{../40.Image/13. Exo-Les parallélogrammes 9.png}
\end{center}

\begin{enumerate}[font=\color{exer1}]
\item Quelle est la longueur $AC$ ? Pourquoi ?
\item Quelle est la longueur $BD$ ? Pourquoi ?
\end{enumerate}
\end{exer1}
\end{minipage}
\hspace{2mm}
\begin{minipage}[t]{8.5cm}
\begin{exer1}
On considère le losange $EFGH$.

\begin{center}
\includegraphics[scale=1]{../40.Image/13. Exo-Les parallélogrammes 10.png}
\end{center}

\begin{enumerate}[font=\color{exer1}]
\item Quelle est la mesure de l'angle $\widehat{EFG}$ ? Pourquoi ?
\item Justifie que les droites $(HF)$ et $(EG)$ sont perpendiculaire.
\end{enumerate}
\end{exer1}
\end{minipage}



\end{document}